% Generated draft: requires mathematical and source review.
\documentclass[11pt]{article}
\usepackage[margin=25mm]{geometry}
\usepackage{fontspec}
\setmainfont{Latin Modern Roman}
\usepackage{amsmath,amssymb,mathtools,unicode-math}
\setmathfont{Latin Modern Math}
\usepackage{enumitem,needspace,xurl,hyperref}
\setlength{\parindent}{0pt}
\setlength{\parskip}{5pt}
\setlength{\emergencystretch}{3em}
\newcommand{\N}{\mathbb{N}}
\newcommand{\Z}{\mathbb{Z}}
\newcommand{\Q}{\mathbb{Q}}
\newcommand{\R}{\mathbb{R}}
\newcommand{\C}{\mathbb{C}}
\newcommand{\F}{\mathbb{F}}
\newcommand{\A}{\mathbb{A}}
\newcommand{\PP}{\mathbb{P}}
\newcommand{\E}{\mathbb{E}}
\newcommand{\eps}{\varepsilon}
\newcommand{\dd}{\,\mathrm{d}}
\begin{document}
\section*{UnsolvedMath ID: 20002478}
\textbf{Source problem code:} catalogue entry
\subsection{Definitions and mathematical statement}
This one registry record merges five consecutive AIM questions, Problems 16--20; all five are retained. Use a smooth autonomous Tonelli Lagrangian on a compact connected manifold \(M\), with fiberwise positive-definite Hessian and superlinear growth, and its dual Hamiltonian \(H:T^*M\to\R\). The source writes \(TM\) for a Hamiltonian in one place; cotangent variables and the canonical symplectic form are intended. Let \(c_u\) be the Ma\~n\'e critical value of the lifted Lagrangian on the universal cover, equivalently the least action shift making the action nonnegative on all closed curves there.

Problem 16 asks whether there are \(M,L\) and a nonempty energy level \(\Sigma_h=\{H=h\}\), with \(h\ne c_u\), containing neither a periodic Hamiltonian orbit nor an equilibrium. The source's statement that only an example at \(c_u\) was known is historical, not a current literature conclusion.

For Problem 17 put \(e_0=\max_{x\in M}\min_{p\in T_x^*M}H(x,p)\), equivalently the least energy whose level projects onto all of \(M\). A regular level is of positive contact type here if there exists a one-form \(\lambda\) on \(\Sigma_h\) with
\[
 d\lambda=\omega_{\mathrm{can}}|_{\Sigma_h},\qquad\lambda(X_H)>0.
\]
This means equality after restriction, not equality of forms on different manifolds; it does not require a primitive on all of \(T^*M\). The source asks whether \(M\ne T^2\) forces every level with \(e_0<h<c_u\) to fail this contact-type condition. Any additional intended dimension or regularity restrictions require source review.

Problem 18 instead takes noncompact \(M\) and a convex superlinear Lagrangian with appropriate, unspecified completeness conditions at infinity. It asks to understand the geometry of a compactification by an ``extended Aubry set'' whose boundary points correspond to Busemann viscosity solutions of the Hamilton--Jacobi equation. The exact compactification topology and equivalence on solutions are not given; the source cites \emph{Calculus of Variations} 13 (2001), 427--458.

Problem 19 takes an autonomous Tonelli Lagrangian on \(T^n\), its dual \(H\), and \(c_0=\alpha(0)\). For the periodic lift \(\widetilde H\) to \(T^*\R^n=\R^{2n}\), it asks whether
\[
 \{(x,p):\widetilde H(x,p)<h\}
\]
has finite symplectic capacity for every \(h<c_0\). The capacity is not specified. The base lift is important: the domain is not a relatively compact subset of the torus cotangent bundle. The value \(c_0\) is not silently identified with \(c_u\).

Problem 20 assumes an autonomous Tonelli \(H\) on \(T^*T^n\), full projection of \(\mathcal A_0\), and uniqueness up to constants of a critical viscosity solution \(u\) of \(H(x,Du)=\alpha(0)\). For the source's sign convention
\[
 H_c(x,p)=H(x,p-c),
\]
let \(v_c\) be a critical viscosity solution and \(\mathcal A_c\) the corresponding Aubry set. Write its critical value as \(\bar\alpha(c)\); this is the source's \(\alpha(c)\), and equals the standard plus-shift \(\alpha\)-function evaluated at \(-c\). The question is
\[
 \lim_{|c|\to0}\|Dv_c-Du\|_{\operatorname{Lip}(\pi(\mathcal A_c))}=0.
\]
On the flat torus, one possible full norm is the supremum plus the Lipschitz seminorm of the vector field on that subset; the source does not state whether it means this full norm or only a seminorm. Gradients of critical solutions on Aubry projections use the usual weak-KAM differentiability there.

The aid's exact special family for Problem 20 is
\[
 H(x,p)=a+K(p-Du(x)),
\]
where \(u\) is smooth periodic and \(K\) is a smooth superlinear momentum Hamiltonian with positive-definite Hessian. It reports
\[
 \bar\alpha(c)=a+K(-c),\qquad v_c=u+\text{constant},\qquad
 \mathcal A_c=\operatorname{graph}(Du)
\]
for the shifted Hamiltonian \(H_c\). Thus the full Lipschitz gradient difference is identically zero for every \(c\), not only in the limit.

The separate quantitative bound assumes a uniformly convex equilibrium critical graph:
\[
 H(x,Du(x))=a,\quad H_p(x,Du(x))=0,\qquad
 mI\preceq H_{pp}(x,p)\preceq MI
\]
with \(0<m\le M\) on the relevant momentum segments. The aid asserts
\[
 \sup_{x\in\pi(\mathcal A_c)}|Dv_c(x)-Du(x)|
 \le(1+\sqrt{M/m})|c|.
\]
This controls a supremum norm, not the Lipschitz seminorm in the general Problem 20. The partial results do not solve Problems 16--19 or the full general Lipschitz question; no new status assessment for those questions is supplied.
\subsection{Short English statement}
This record contains five distinct Ma\~n\'e-critical and Aubry questions. The aid's partial results concern only the fifth: exact momentum-gauge Hamiltonians give zero Lipschitz gradient error, and a more general equilibrium graph gives a uniform gradient bound that does not control the full Lipschitz seminorm.
\subsection{Sources}
\begin{itemize}
\item Local registry, permanent UnsolvedMath ID \(20002478\), exact public code \texttt{catalogue entry}; canonical location \texttt{aim-pdes-notes.json}, \texttt{notes[89]}.
\item AIM, \emph{New connections between dynamical systems and PDE's}, 2003 workshop list, Problems 16--20 merged as one record; supplied reference \url{https://aimath.org/WWN/dynpde/dynpde.pdf}. No PDF was accessed here. No proposer, original remarks, or separate literature field is supplied.
\item Problem 18 cites \emph{Calculus of Variations} \textbf{13} (2001), 427--458, without author, title, or theorem number. The historical ``only example known'' assertion belongs to Problem 16's original wording.
\item Supplied research attempt 1: \texttt{partial\_result}, run and registry status \texttt{partially\_solved}, rigor label \texttt{proved}, low novelty confidence; classification \texttt{PARTIAL-PROGRESS}, created and updated 14 August 2026. Its exact-gauge family and equilibrium-graph estimate have no primary bibliographic citation in the record. The aid is machine generated, not peer reviewed, and requires independent verification before citation.
\end{itemize}
\end{document}
